\subsection{Characteristic two, the universal support, and all jet lengths}
\label{ir:sec:modular}

\subsubsection{A characteristic-two homology theorem over arbitrary algebras}
Assume now that the coefficient ring $S$ has characteristic two.
Write $N=c-1=c+1$. Then $N^2=0$, and both periodic cohomology
groups are the same module $\ker N/\operatorname{im} N$. On a fixed-point module
both periodic differentials vanish. The comparison
\eqref{ir:eq:tate-comparison} therefore has a particularly simple
consequence.

\begin{theorem}[Koszul model in characteristic two]
\label{ir:thm:char2-koszul}
For every commutative characteristic-two ring $S$ and every choice
of its weights,
\begin{equation}\label{ir:eq:char2-homology}
 \frac{\ker(N:\mathcal D_q(S,a)\to\mathcal D_q(S,a))}{\operatorname{im} N}
 \simeq\bigoplus_{j=0}^{r(q)}H_j(K_S(\boldsymbol\delta)),
\end{equation}
where the active sequence is \eqref{ir:eq:active-koszul}.
\end{theorem}
\begin{proof}
Remove the moving-orbit subcomplex as in
Section~\ref{ir:sec:reflection}. On its fixed quotient the vertical
periodic differential is zero. Total cohomology is therefore the
direct sum of the horizontal homology groups, one from each of the
finitely many chain degrees. Apply the chain-homotopy reduction of
Proposition~\ref{ir:prop:fixed-koszul}.
\end{proof}

\subsubsection{The universal obstruction is a reduced complete intersection}
Let $R_2=\mathbb F_2[A_p:p\mid q]$ and let $I_q\subseteq R_2$ be the ideal
\begin{equation}\label{ir:eq:universal-ideal}
 I_q=(A_p-1:\ p>2,\ p\mid q)
       +(A_2(A_2-1):\ 4\mid q),
\end{equation}
where the second summand is omitted unless indicated. Its generators
form a regular sequence. Indeed the odd-prime entries are independent
linear variables shifted by one; after removing them, the possible
last entry is a nonzero monic polynomial in the still-free variable
$A_2$. At each step multiplication by that entry is injective in the
quotient by the earlier ones.

The Koszul complex on a regular sequence resolves its quotient. This
can be seen inductively here without a general regular-sequence
criterion: tensor a free resolution of the previous quotient with
$[R_2\xrightarrow{\delta}R_2]$. Its homology consists of the kernel
and cokernel of multiplication by $\delta$ on that quotient. The
kernel is zero, and the cokernel is the next quotient.

\begin{corollary}[Universal characteristic-two support]
\label{ir:cor:universal-support}
Over $R_2$, the reflection homology is a cyclic module:
\begin{equation}\label{ir:eq:universal-homology}
 \ker N/\operatorname{im} N\simeq R_2/I_q.
\end{equation}
Its zeroth Fitting ideal is exactly $I_q$. The closed locus defined
by $I_q$ is reduced. For odd $q>1$ it is the single point with every
weight equal to one. For $4\mid q$ it has two points, with all
odd-prime weights equal to one and $A_2=0$ or $1$. For $q$ twice an
odd number, $A_2$ is free and the locus is an affine line.
\end{corollary}
\begin{proof}
The positive Koszul homology vanishes by the preceding argument;
apply Theorem~\ref{ir:thm:char2-koszul}. A cyclic module $R_2/I_q$
has zeroth Fitting ideal $I_q$ by its one-generator presentation.
After the odd linear equations, the only possible further equation
is $A_2(A_2-1)$, whose two factors are relatively prime. Its quotient
is a product of two fields, or of two copies of the remaining
coefficient ring if one extends scalars. The stated descriptions
and reducedness follow.
\end{proof}

The cyclic universal answer does \emph{not} imply that every fiber
of the reflection homology is one-dimensional. Homology need not
commute with a nonflat specialization. At an exceptional field-valued
point the differentials of the specialized Koszul complex all vanish,
so its homology has total dimension $2^{r(q)}$. This is the precise
source of the larger fiber defect.

There is also an integral universal formulation. With the periodic
convention \eqref{ir:eq:tate-definition}, the same fixed-complex
calculation, followed by reduction modulo two, gives
\begin{equation}\label{ir:eq:universal-integral-tate}
 \widehat H^{2j}(\mathcal D_q)\simeq R_q/(2,\boldsymbol\delta),\qquad
 \widehat H^{2j+1}(\mathcal D_q)=0.
\end{equation}
Here $\mathcal D_q$ is still over the polynomial ring, not an integer weight
fiber. The regularity argument is in $R_q/2R_q$. Odd periodic
cohomology appearing after integer specialization is another
nonflat-base-change effect, not a contradiction to
\eqref{ir:eq:universal-integral-tate}.

\subsubsection{The exact rank locus over every field}
For a field $k$ of characteristic two, a nonzero active scalar is a
unit and makes the Koszul complex contractible. If all active scalars
vanish, its differentials vanish and its total homology dimension is
$2^{r(q)}$. To pass from homology to the desired coinvariants, note
that an endomorphism $N$ with $N^2=0$ on a $d$-dimensional vector
space satisfies
\begin{equation}\label{ir:eq:nilpotent-dimension}
 \dim(\ker N/\operatorname{im} N)=d-2\rank N,
 \qquad
 \dim\operatorname{coker} N=\frac{d+\dim(\ker N/\operatorname{im} N)}2.
\end{equation}

\begin{theorem}[All-field reflection dimension]
\label{ir:thm:field-locus}
Let $q>2$, let $k$ be a field, and use arbitrary weights $a_p\in k$.
If $\operatorname{char}k\ne2$, both sign coinvariants have dimension
$n(q)$. If $\operatorname{char}k=2$, the signs coincide and
\begin{equation}\label{ir:eq:field-locus}
 \dim_k\mathcal D_q(k,a)/(c-1)\mathcal D_q(k,a)
 =n(q)+h(q)\mathbf 1_{\mathcal E_q(a)},
\end{equation}
where $\mathcal E_q(a)$ means
\[
 a_p=1\quad\text{for every odd }p\mid q,
 \qquad a_2\in\{0,1\}\quad\text{if }4\mid q.
\]
There is no condition on $a_2$ when $v_2(q)=1$.
\end{theorem}
\begin{proof}
For characteristic different from two use
Lemma~\ref{ir:lem:balanced}. In characteristic two use the Koszul
calculation and \eqref{ir:eq:nilpotent-dimension} with
$d=\varphi(q)$. Over a field, $a_2(a_2-1)=0$ is equivalent to
$a_2=0$ or $1$.
\end{proof}

Over $\mathbb F_2$, the condition at $2$ is automatic. Over $\mathbb F_4$ it is
not: a weight satisfying $a_2^2+a_2+1=0$ is neither zero nor one,
and the reflection defect disappears when $4\mid q$. Thus an
experiment restricted to binary weights would miss part of the
geometry. The exact regression suite includes all four field
weights at selected levels to test this distinction.

\subsubsection{A complete one-parameter Smith law}
Let $k$ have characteristic two, let $O=k[[t]]$, and choose power
series weights $a_p(t)$. The module $\mathcal D_q(O,a)$ is free of rank
$2n=\varphi(q)$. Form the active sequence $\boldsymbol\delta(t)$
as before, and define its contact order by
\begin{equation}\label{ir:eq:contact-order}
 d=\min_i v_t(\delta_i(t)),\qquad v_t(0)=\infty.
\end{equation}
The word ``contact'' refers to the pullback of the explicit ideal
$I_q$ along the chosen formal curve. It is not an approximation to a
root of a transcendental function.

Each active prime contributes a factor at least two to $\varphi(q)$,
so $\varphi(q)\ge2^{r(q)}$ and $n\ge h$. Thus all block sizes in
the following statement are nonnegative.

\begin{theorem}[Contact-order Smith law]\label{ir:thm:smith-jets}
Let $q>2$, $n=n(q)$, $h=h(q)$, and use the above characteristic-two
power series weights. If $d<\infty$, the Smith diagonal of
$N=c-1$ over $O$ is, up to units,
\begin{equation}\label{ir:eq:smith-contact}
 \operatorname{diag}\bigl(
    \underbrace{1,\ldots,1}_{n-h},
    \underbrace{t^d,\ldots,t^d}_{h},
    \underbrace{0,\ldots,0}_{n}\bigr).
\end{equation}
When $d=0$, the first two blocks together are $n$ unit entries.
When $d=\infty$, replace the $h$ entries $t^d$ by zero; there are
then $n-h$ unit entries and $n+h$ zeros.

For every integer $L\ge1$, put $O_L=k[t]/(t^L)$. The complete jet
cokernel formula is
\begin{equation}\label{ir:eq:jet-length}
 \dim_k\frac{\mathcal D_q(O_L,a)}{(c-1)\mathcal D_q(O_L,a)}
 =nL+h\min(d,L),
\end{equation}
with $\min(\infty,L)=L$.
\end{theorem}

\begin{proof}
First work over $O_L$ and put $d_L=\min(d,L)$. Elementary invertible
operations on the active sequence identify its Koszul complex with
that on $(t^{d_L},0,\ldots,0)$, with the convention that $t^L=0$.
To see this explicitly when $d_L<L$, choose an entry with least
valuation, multiply it by a unit to make it $t^{d_L}$, and subtract
its multiples from all other entries. Every other entry is divisible
by it in $O_L$. If $d_L=L$, all entries already vanish. Invertible
changes of the sequence induce invertible changes on its exterior
algebra, hence isomorphisms of Koszul complexes.

The kernel and cokernel of multiplication by $t^{d_L}$ on $O_L$
each have $k$-dimension $d_L$. Tensoring with $r-1$ zero-differential
two-term complexes multiplies total homology dimension by $2^{r-1}$.
The right-hand side of \eqref{ir:eq:char2-homology} therefore has
$k$-dimension $2^r d_L$. The underlying vector space of
$\mathcal D_q(O_L,a)$ has dimension $2nL$, by the integral basis theorem.
Applying \eqref{ir:eq:nilpotent-dimension} proves
\eqref{ir:eq:jet-length}.

It remains to show that these length formulas determine all Smith
exponents, rather than just their sum. A matrix over the discrete
valuation ring $O$ has nonzero invariant factors $t^{e_i}$ up to
units, together with some zero factors. If there are $z$ zero
factors, its cokernel modulo $t^L$ has length
\[
 zL+\sum_i\min(e_i,L).
\]
The first difference in $L$ is $z+\#\{i:e_i\ge L\}$ for $L\ge1$.
For finite $d$, the formula already proved has first difference
$n+h$ for $1\le L\le d$, and $n$ for $L>d$. Thus exactly $n$
factors are zero and exactly $h$ nonunit factors have exponent $d$;
the remaining $n-h$ factors are units. For $d=0$ there are no
positive exponents. For $d=\infty$, the first difference is always
$n+h$, so there are $n+h$ zero factors and no positive finite
exponents. This proves every case of \eqref{ir:eq:smith-contact}.
\end{proof}

\subsubsection{Examples of the jet law}
At level $15$, take
\[
 a_3(t)=1+t^2,\qquad a_5(t)=1+t^5.
\]
Here $(n,r,h)=(4,2,2)$ and $d=2$. The reflection matrix has Smith
entries $1,1,t^2,t^2,0,0,0,0$, and its jet cokernel dimensions at
$L=1,2,3,4$ are $6,12,16,20$. A rank calculation at $t=0$ alone
would detect the first number, but not the exponent two.

At level $12$, take $a_2(t)=t^3$ and $a_3(t)=1+t^5$.
Then $a_2(a_2-1)=t^3(1+t^3)$, so $d=3$. Since $(n,h)=(2,2)$,
there are no unit Smith factors and the diagonal is
$t^3,t^3,0,0$. Every jet of length at most three has trivial
reflection matrix up to equivalence; later lengths detect the
nonzero factors.

If the entire curve is contained in the exceptional locus, then
$d=\infty$ and the dimension is $(n+h)L$. If any active scalar
has a nonzero constant term, $d=0$ and the dimension is $nL$.
These are exact boundary cases of the same formula. In contrast,
the \emph{unreflected} distribution module always has dimension
$\varphi(q)L$ at length $L$ and never exhibits this jump.
